\section{Reductions and a finite free complex}\label{sec:reductions}

For maps of positive residue characteristic, our aim is to obtain
a factorization
$R\to T\twoheadrightarrow S$ with $e(T)=e(R)$ and a short free
complex over $T$ supported at its closed point. The first reduction
applies in every characteristic and handles arbitrary dimensions
and residue fields. The construction
that follows also controls the generic local rings of $T$; these
will determine the coefficients in the final multiplicity sums.

\subsection{Completion, components, and residue fields}

We first make reductions which do not impose any restriction on the
characteristic or on the original residue fields.  We use the following
localization theorem for Hilbert--Samuel multiplicity: if $(B,\mathfrak b)$
is an excellent local ring and $\mathfrak q\in\Spec B$ satisfies
\[
  \dim(B/\mathfrak q)+\dim B_{\mathfrak q}=\dim B,
\]
then $e(B_{\mathfrak q})\leq e(B)$.  This is Nagata's localization theorem;
the stated form is recalled in \cite[Theorem~2.1]{Ma}.  It does not require
$B$ to be equidimensional.

\begin{lemma}\label{lem:residue-extension}
For every complete Noetherian local ring $(B,\mathfrak b,k)$ and every
algebraic closure $\overline{k}$ of $k$, there is a complete Noetherian
local ring $B'$ and a flat local map $B\to B'$ such that
\[
  \mathfrak bB'=\mathfrak b_{B'},\qquad
  B'/\mathfrak b_{B'}\cong\overline{k},\qquad e(B')=e(B).
\]
\end{lemma}

\begin{proof}
In equal characteristic, choose a coefficient field and a presentation
$B=k[[X_1,\ldots,X_t]]/I$, and set
\[
  B'=\overline{k}[[X_1,\ldots,X_t]]/
             I\overline{k}[[X_1,\ldots,X_t]].
\]
The map between the power series rings is flat.  For example, this
follows from the local flatness criterion modulo $(X_1,\ldots,X_t)$:
the variables form regular sequences in both rings, and the resulting
map of coefficient fields is flat.  The quotient map is therefore flat
as well and has the asserted maximal ideal and residue field.

In the remaining case the residue characteristic is $p>0$, and the
Cohen Structure Theorem~\cite[Tag~032A]{Stacks} gives a presentation
$B=V[[X_1,\ldots,X_t]]/I$, where $V$ is a complete discrete
valuation ring with uniformizer $p$ and residue field $k$.  We explain
why $V$ admits a complete discrete valuation extension $V'$ with the
same uniformizer and residue field $\overline{k}$, even when $k$ is
imperfect.  Given a simple algebraic residue extension with irreducible
minimal polynomial $\overline f$, lift $\overline f$ to a monic
polynomial $f\in V[X]$.  The algebra $V[X]/(f)$ is finite free over $V$,
and its quotient by $p$ is the desired residue field.  Every maximal
ideal lies over $(p)$, so this algebra is local with maximal ideal
generated by $p$.  The element $p$ is a nonzerodivisor, and the algebra
has dimension one.  It is consequently a discrete valuation ring.
This argument does not require $\overline f$ to be separable.

Repeat this construction along a tower of simple algebraic extensions
whose union is $\overline{k}$.  At limit stages take unions.  Every
nonzero element of such a union is $p^a$ times a unit for an integer
$a\geq0$; hence every nonzero ideal is generated by an element of least
valuation.  The union is again a discrete valuation ring with
uniformizer $p$.  Completing the final union gives the required $V'$.
The extension $V\to V'$ is flat because $V'$ is torsion-free over $V$.
Now use
\[
  B'=V'[[X_1,\ldots,X_t]]/IV'[[X_1,\ldots,X_t]].
\]
Flatness of the power series extension follows from the same local
flatness criterion, this time with coefficient map $V\to V'$.

Finally, for any flat local map with extended maximal ideal, a
composition series of a finite-length module remains a composition
series after tensoring: each residue-field factor becomes the target
residue field.  In particular
\[
  \length_{B'}(B'/\mathfrak b^N B')
       =\length_B(B/\mathfrak b^N)\quad(N\geq1).
\]
The two Hilbert--Samuel functions, and therefore their degrees and
multiplicities, agree.
\end{proof}

\begin{lemma}\label{lem:reduction}
To prove $e(R)\leq e(S)$ for all flat local maps of nonzero Noetherian
local rings, it suffices to prove it when $R$ and $S$ are complete,
$R$ is a domain, $\dim R=\dim S>0$, and the residue field of $S$ is
algebraically closed.  The reduction preserves equal characteristic
$p$ when both original rings have characteristic $p$.
\end{lemma}

\begin{proof}
Completion preserves dimension and Hilbert--Samuel multiplicity, and
a flat local map induces a flat map on completions.  We can thus first
assume that $R$ and $S$ are complete.  Put $d=\dim R$ and
$f=\dim(S/\mathfrak mS)$.  The dimension formula for a flat local map
gives $\dim S=d+f$.  Choose $Q$ minimal over $\mathfrak mS$ with
$\dim(S/Q)=f$.  Going-down gives $\dim S_Q\geq d$, whereas
\[
  \dim S\geq\dim(S/Q)+\dim S_Q\geq f+d=\dim S.
\]
All these inequalities are equalities.  Since $S$ is excellent,
Nagata's theorem gives $e(S_Q)\leq e(S)$.  The map $R\to S_Q$ is flat
local, its closed fiber has dimension zero, and $\dim S_Q=d$.
Replacing $S_Q$ by its completion preserves these properties.  Thus a
proof for this new target proves the original inequality.

Choose a finite prime filtration of the $R$-module $R$, with factors
$R/\mathfrak p_j$, allowing repetitions.  Flatness gives a filtration
of $S$ with factors $S/\mathfrak p_jS$.  The flat local map
\[
  R/\mathfrak p_j\longrightarrow S/\mathfrak p_jS
\]
has the same zero-dimensional closed fiber, so its source and target
have equal dimension.  The additivity formula for multiplicity yields
\[
 e(R)=\sum_{\dim(R/\mathfrak p_j)=d}e(R/\mathfrak p_j),\qquad
 e(S)=\sum_{\dim(R/\mathfrak p_j)=d}e(S/\mathfrak p_jS).
\]
It therefore suffices to prove the inequality for the indicated domain
sources.  Their quotients and the corresponding targets remain
complete.  If $d=0$, each such source is a field and has multiplicity
one, while a nonzero zero-dimensional local target has multiplicity
equal to its positive length.  This case is settled.

For $d>0$, apply Lemma~\ref{lem:residue-extension} to the target.
The resulting composite remains flat local and has the same target
dimension and multiplicity.  This proves the reduction.  The argument
also appears in \cite[Lemma~2.2]{Ma}.
\end{proof}

\subsection{A factorization with a perfect kernel}

The passage from a flat local map to a short free complex follows the
construction used by Iyengar--Ma--Walker
\cite[Proposition~8.3]{IMW2022}. We construct one auxiliary ring and one
complex for both cases of positive residue characteristic, while also
controlling the local rings needed to sum the component estimates. A
\emph{perfect ideal} $J$ of a Noetherian local ring $T$ means an ideal
for which $T/J$ has finite projective dimension and
\[
  \operatorname{pd}_T(T/J)=\operatorname{grade}_T J.
\]

\begin{proposition}\label{prop:setup}
Let $(R,\mathfrak m)\to(S,\mathfrak n)$ be a flat local map of complete
Noetherian local rings.  Suppose $R$ is a domain,
$\dim R=\dim S=d>0$, the residue characteristic is $p>0$, and the
residue field $L$ of $S$ is algebraically closed.  There is a factorization
\[
  R\longrightarrow (T,\mathfrak t)\longrightarrow S=T/J
\]
with the following properties.  Write
$b=\dim(T/\mathfrak mT)$ and $h=d+b$.
\begin{enumerate}
\item[(i)] The ring $T$ is complete, its residue field is $L$, and
      $R\to T$ is flat with regular closed fiber.  Moreover
      $\dim T=h$ and $e(T)=e(R)$.
\item[(ii)] Every irreducible component of $T$ has dimension $h$.
      The ideal $J$ is perfect of grade $b$, and $S$ has a minimal
      finite free $T$-resolution $P_\bullet$ in degrees $0,\ldots,b$
      with $P_0=T$.
\item[(iii)] Choose lifts $x_1,\ldots,x_d\in\mathfrak t$ of a minimal
      reduction of $\mathfrak n$.  Then
      \[
         F_\bullet=P_\bullet\otimes_T K_\bullet(x_1,\ldots,x_d;T)
      \]
      is a finite free complex in degrees $0,\ldots,h$, has $F_0=T$,
      and is contractible locally on
      $\Spec T\setminus\{\mathfrak t\}$.  Its differential entries
      lie in $\mathfrak t$, and its ranks have alternating sum zero.
\item[(iv)] If $Q\in\Spec T$ contracts to zero in $R$, then $T_Q$ is
      a complete intersection.  In mixed characteristic, $T_Q$ is
      regular.  Every prime of $T$ which is minimal over $J$ contracts
      to zero in $R$.  In mixed characteristic $T$ is generically
      reduced; a prime $Q$ minimal over $J$ contains a unique minimal
      prime of $T$.
\end{enumerate}
\end{proposition}

Parts (i)--(iii) preserve the source multiplicity and produce the complex
to which Theorem~\ref{thm:complex} will apply. Part (iv) supplies the
characteristic-dependent comparisons at primes $Q$ minimal over $J$.
In characteristic $p$, its complete-intersection assertion will identify
$e_{\mathrm{HK}}(JT_Q,T_Q)$ with $\length_{T_Q}(T_Q/JT_Q)$ in
Lemma~\ref{lem:ma-hk}. In mixed characteristic, regularity puts each
such prime on exactly one component of $T$, so the later summation
counts it only once.

\begin{proof}
\emph{Constructing $T$.}
Let $k$ be the residue field of $R$.  Choose a complete regular
normalization $A\subset R$ with residue field $k$, such that $R$ is
finite over $A$.  In equal characteristic, take a coefficient field
$k\subset R$ and a system of parameters $a_1,\ldots,a_d$, giving
$A=k[[a_1,\ldots,a_d]]$.  In mixed characteristic, choose a coefficient
Cohen ring $V\subset R$ and a system of parameters
$p,a_2,\ldots,a_d$, giving $A=V[[a_2,\ldots,a_d]]$.  The Cohen
Structure Theorem and completeness give finiteness in both cases;
the maps from these power series rings are injective by dimension.
No generic separability is required here.

Apply the Cohen Factorization Theorem of
Avramov--Foxby--Herzog~\cite{AFH} to $A\to S$:
\[
  A\longrightarrow (U,\mathfrak u)\twoheadrightarrow S,
\]
where $U$ is complete, $A\to U$ is flat local, and its closed fiber is
regular; see \cite[Theorem~3.1]{Ma}.  Regularity ascends along this map
because both $A$ and the closed fiber are regular
\cite[Tag~031E]{Stacks}. Thus $U$ is a
regular local ring, with residue field $L$.

Set $T=R\otimes_A U$.  This is finite over $U$, and multiplication of
the two given maps to $S$ induces a surjection $T\to S$.  Its closed
fiber over $U$ is
\[
  T/\mathfrak uT=(R/\mathfrak m_A R)\otimes_k L.
\]
The ring $R/\mathfrak m_A R$ is local Artinian with residue field $k$.
Its radical is nilpotent, so the displayed fiber is local Artinian
with residue field $L$.  A finite algebra over a complete local ring
is a product of complete local rings; this fiber shows that there is
only one factor.  Hence $T$ is complete local with residue field $L$.
Base change shows that $R\to T$ is flat and
\[
  T/\mathfrak mT\cong U/\mathfrak m_A U.
\]
Its closed fiber is therefore regular of dimension $b$, and the flat
dimension formula gives $\dim T=d+b=h=\dim U$.

\begin{figure}[ht]
\centering
\[
\begin{array}{ccccc}
 A&\longrightarrow&U&&\\
 \big\downarrow&&\big\downarrow&&\\
 R&\longrightarrow&T=R\otimes_A U&\twoheadrightarrow&S .
\end{array}
\]
\caption{The construction of the auxiliary factorization. The square
is a tensor-product square; $R$ is finite over $A$, and $A\to U$
is flat local with regular closed fiber. The surjection to $S$
comes from its given $R$- and $U$-algebra structures.}
\label{fig:factorization}
\end{figure}
Figure~\ref{fig:factorization} records the two operations in the
construction: the finite normalization $A\subset R$ and the Cohen
factorization of $A\to S$.

\smallskip
\noindent\emph{Preserving multiplicity.}
To check the multiplicity assertion, lift a regular system of
parameters of $T/\mathfrak mT$ to $w_1,\ldots,w_b\in\mathfrak t$.
Lifting regular sequences under a flat local map
\cite[Tag~00MG]{Stacks} shows that $w$ is a
$T$-regular sequence and $T/(w)$ is flat over $R$.  Completeness
defines the evaluation map
\[
   R[[W_1,\ldots,W_b]]\longrightarrow T,\qquad W_i\longmapsto w_i.
\]
The local flatness criterion modulo the variables
\cite[Tag~00ML]{Stacks} makes this map flat.
Its extended maximal ideal is $\mathfrak t$, because the images of
$w$ generate the maximal ideal of the closed fiber.  As in the last
paragraph of Lemma~\ref{lem:residue-extension}, the Hilbert--Samuel
functions are unchanged under a flat local map with extended maximal
ideal.  Adjoining formal variables preserves multiplicity, as follows
also from the associated graded polynomial extension.  Therefore
\[
  e(T)=e(R[[W_1,\ldots,W_b]])=e(R).
\]

\smallskip
\noindent\emph{The resolution and the perfect kernel.}
Start with a minimal free $T$-resolution of $S$, possibly infinite.
Each syzygy is flat over $R$:
inductively, the kernel of a surjection between $R$-flat modules with
flat quotient is flat.  Tensoring with $k$ over $R$ consequently
preserves exactness and yields a minimal free resolution of
$S/\mathfrak mS$ over the regular local ring $T/\mathfrak mT$.
That ring has global dimension $b$.  Minimality, or Nakayama's Lemma,
therefore forces the original resolution to vanish in degrees greater
than $b$.  Thus $\operatorname{pd}_T S\leq b$.

The closed fiber $S/\mathfrak mS$ is zero-dimensional, so the image of
$J$ in $T/\mathfrak mT$ is primary for its maximal ideal.  Choose in
this image a system of $b$ parameters and lift them to elements of
$J$.  Since the fiber is regular, this is a regular sequence on the
fiber; flat regular-sequence lifting makes its lifts a $T$-regular
sequence.  Consequently
\[
 b\leq\operatorname{grade}_T J
   \leq\operatorname{pd}_T S\leq b.
\]
This proves perfection and identifies the length of the minimal
resolution.  The assertions include $b=0$, when the regular sequence
is empty and the resolution has only its degree-zero term.

\smallskip
\noindent\emph{Components and the supported complex.}
Because $R$ is a finite torsion-free $A$-module, it embeds in a finite
free $A$-module.  Flatness gives an embedding of $T$ in a finite free
$U$-module.  The contractions to $U$ of the associated primes of $T$
are thus zero.  In particular every minimal prime $P$ of $T$ contracts
to zero, and $T/P$ is finite integral over $U$.  It follows that
$\dim(T/P)=\dim U=h$.

The field $L$ is infinite, so a minimal reduction of $\mathfrak n$
generated by $d$ elements exists.  Lift its generators to $x$ in $T$.
The ideal $J+(x)$ is $\mathfrak t$-primary.  At a prime outside its
vanishing set, either the localized resolution resolves zero and is
split exact, or one of the $x_i$ is a unit and its Koszul complex is
contractible.  Hence $F_\bullet$ is contractible off $\mathfrak t$.
Minimality of the resolution and $x_i\in\mathfrak t$ put all its
differential entries in $\mathfrak t$.  Its degree-zero term is $T$,
and
\[
  \sum_i(-1)^i\rank F_i
    =\left(\sum_i(-1)^i\rank P_i\right)(1-1)^d=0,
\]
where $d>0$ is used.

\smallskip
\noindent\emph{Local rings over the generic point of $R$.}
Finally put $K=\operatorname{Frac}(A)$ and
$K_R=\operatorname{Frac}(R)$.  Localizing at $A\setminus\{0\}$ gives
\[
 T\otimes_A K
   \cong (U\otimes_A K)\otimes_K K_R.
\]
Here $U\otimes_A K$ is regular, and $K_R/K$ is a finite field
extension.  Express this field extension as a finite tower of simple
algebraic extensions.  After base change each step is defined by a
monic polynomial, which is a nonzerodivisor even if the intermediate
ring is not reduced.  The resulting local rings are quotients of
regular local rings by regular sequences, hence complete
intersections.  Every $Q$ contracting to zero in $R$ belongs to this
localization, proving the first assertion of (iv).  In mixed
characteristic the fields have characteristic zero, so $K_R/K$ is
separable.  The displayed extension is then finite etale, and its
local rings are regular.

Every minimal prime of $S$ contracts to a minimal prime of $R$ by
flat going-down; that prime is zero.  This applies to the primes of
$T$ minimal over $J$.  Moreover every minimal prime of $T$ avoids
$A\setminus\{0\}$ by the preceding torsion-freeness argument, so in
mixed characteristic its localization is a zero-dimensional regular
local ring, namely a field.  Thus $T$ is generically reduced.  At a
prime $Q$ minimal over $J$, the regular local ring $T_Q$ is a domain,
which also shows that $Q$ contains exactly one minimal prime of $T$.
\end{proof}

\begin{remark}\label{rem:setup-dimensions}
Since $T$ is complete and equidimensional, it is catenary and
\[
  \dim T_Q+\dim(T/Q)=h\qquad(Q\in\Spec T).
\]
In particular a prime $Q$ containing $J$ with $\dim(T/Q)=d$ is minimal
over $J$ and has $\dim T_Q=b$.  In mixed characteristic its
localization is regular by Proposition~\ref{prop:setup}.
\end{remark}
